\section{数据并行（Data Parallelism）}

\subsection{基本原理}

数据并行是最简单也是最基础的分布式训练方法。其核心思想：每块 GPU 持有模型的完整副本，但处理不同的数据子集。

\begin{figure}[H]
\centering
\includegraphics[width=0.95\textwidth]{slides-images/slide-013.jpg}
\caption{数据并行的数学基础：梯度的线性性使得可以先在每个 GPU 上独立计算局部梯度，再进行全局平均\protect\footnotemark}
\end{figure}
\footnotetext{Slides 第13页。}

数学上，设总损失为各样本损失的平均：

$$L = \frac{1}{M \cdot N} \sum_{j=1}^{M} \sum_{i=1}^{N} \ell(x_{ij}, W)$$

由于梯度是线性算子，可以重新排列为：

$$\nabla_W L = \frac{1}{M} \sum_{j=1}^{M} \underbrace{\left[\frac{1}{N} \sum_{i=1}^{N} \nabla_W \ell(x_{ij}, W)\right]}_{\text{GPU } j \text{ 独立计算}}$$

\begin{itemize}
    \item $M$：GPU 数量
    \item $N$：每块 GPU 上的 mini-batch 大小
    \item 蓝色括号内的部分可以在每块 GPU 上\textbf{完全独立}计算
    \item 外层求和需要\textbf{跨 GPU 通信}（All-Reduce）
\end{itemize}

\subsection{数据并行的执行流程}

\begin{figure}[H]
\centering
\includegraphics[width=0.95\textwidth]{slides-images/slide-014.jpg}
\caption{数据并行的六步执行流程\protect\footnotemark}
\end{figure}
\footnotetext{Slides 第14页。}

\begin{enumerate}
    \item 每块 GPU 持有模型权重的\textbf{独立但相同}的副本
    \item 每块 GPU 加载\textbf{不同的} mini-batch 数据
    \item 每块 GPU 独立执行前向传播，计算局部损失
    \item 每块 GPU 独立执行反向传播，计算局部梯度
    \item \textbf{All-Reduce}：所有 GPU 交换并平均梯度
    \item 每块 GPU 独立更新权重（因为起点相同、梯度相同，更新后的权重仍然相同）
\end{enumerate}

\begin{warningbox}{不同 GPU 必须加载不同的数据}
这是数据并行中最容易犯的 bug：所有 GPU 意外加载了相同的 mini-batch。这不会导致程序报错，但会使训练完全无效——你相当于在浪费 $(M-1)$ 块 GPU 的计算资源。
\end{warningbox}

\begin{knowledgebox}{计算与通信的重叠}
反向传播和 All-Reduce 可以\textbf{流水线化执行}：当 GPU 在计算第 $L-1$ 层的梯度时，第 $L$ 层的梯度已经可以开始 All-Reduce。这种重叠隐藏了大部分通信开销。PyTorch 的 \texttt{DistributedDataParallel}（DDP）类自动实现了这一优化。
\end{knowledgebox}

\begin{figure}[H]
\centering
\includegraphics[width=0.95\textwidth]{slides-images/slide-016.jpg}
\caption{计算与通信的流水线化：反向传播（黑色）与 All-Reduce（红色）同时进行\protect\footnotemark}
\end{figure}
\footnotetext{Slides 第16页。}

\subsection{数据并行的瓶颈：内存限制}

数据并行要求每块 GPU 存储模型的完整副本。对于每个模型参数，需要存储：

\begin{itemize}
    \item 权重本身（2 bytes，16-bit）
    \item 梯度（2 bytes）
    \item 优化器状态——Adam 的 $\beta_1$ 和 $\beta_2$（各 2 bytes）
\end{itemize}

即每个参数需要 $\sim$8 bytes。H100 有 80 GB 内存，因此数据并行最多支持约 \textbf{10B 参数}的模型——远远不够。

\subsection{本章小结}

数据并行是分布式训练的基础，数学上完全等价于单 GPU 大 batch 训练。它简单有效，但受限于 GPU 内存——每块 GPU 必须能容纳完整的模型参数、梯度和优化器状态。

%% ========== FSDP ==========

